高维足式机器人控制的核心难点在于：控制器必须在极短的滚动优化周期内，同时处理非线性刚体动力学、混杂接触模式、执行器约束以及与机构几何紧密相关的可行性条件。对于六足机器人而言，这一问题更加突出，因为更多的腿意味着更高的接触冗余、更大的决策空间，以及更多任务空间近似等价但关节空间物理意义不同的解。

本文研究的正是这样一种在六足机器人上真实出现的失效模式。在基于 OCS2 的前向行走闭环中，我们发现原有 SQP 路径可能会收敛到一个局部上“看起来合理”、但在机构分支意义上错误的 tarsus 解。尽管该解在足端任务空间上与目标较为接近，但它落在了左右 tarsus 关节的错误物理解分支上。runtime branch guard 的确可以在求解后拦截这些非法命令，避免直接执行危险动作；然而，这种机制只是在结果层面止损，并没有修复导致问题产生的最优控制问题本身。其直接后果是：branch guard 高频触发、机器人长时间停滞、前向推进几乎消失。

本文的核心观点是：分支一致性必须在求解器内部被编码，而不能仅依赖优化后的 runtime 修补。因此，我们在 OCS2 问题定义中引入硬约束 \emph{TarsusBranchConstraint}，并将求解后端由原来的 SQP 切换为内点法（IPM），使该不等式约束在求解过程中真正参与可行域塑形。这样一来，物理上错误的分支不再只是“代价更高”或“事后被拦截”，而是从求解器有效搜索空间中被直接移除。

本文的主要贡献如下：
\begin{enumerate}
    \item 揭示并定位了一类此前容易被忽视的六足机器人 MPC 失效模式，即 warm-start 驱动的 SQP 在任务空间表面合理的前提下，仍可能收敛到物理上错误的 tarsus 分支。
    \item 提出求解器可见的分支一致性建模方式，通过硬约束 \emph{TarsusBranchConstraint} 将左右 tarsus 的物理解条件直接写入最优控制问题，并说明切换至 IPM 如何恢复真实的不等式约束执行。
    \item 在 Isaac Sim 全闭环实验中给出可复现验证，证明所提方法能够将 branch-guard 触发降为零，恢复稳定前向行走，并在保持亚 5~ms 平均求解时间的前提下满足实时控制要求。
\end{enumerate}

更广义地看，本文表明：多足机器人中的安全性不能仅靠目标函数整形或 runtime 过滤来保证。当机构层面的可行性没有进入优化问题本身时，局部极小值可能会呈现出极强的欺骗性。通过将物理分支有效性显式写入可行域，并采用能够真实处理该不等式的求解器，我们得到的是一个不仅更稳定，而且在科学解释上更清晰的控制系统。
